\chapter{Análisis con MITRE ATT\&CK}
\label{ch:mitre_attack}

Para caracterizar el comportamiento táctico de los adversarios, se utiliza la matriz \textbf{MITRE ATT\&CK for Enterprise}, permitiendo identificar brechas de visibilidad e implementar reglas de detección precisas.

\section{Estudio de TTPs Aplicables a la Superficie de la UNAP}
A continuación, se describen las Tácticas, Técnicas y Procedimientos (TTPs) más relevantes para la infraestructura evaluada:

\begin{itemize}
    \item \textbf{Reconocimiento (TA0043):}
    \begin{itemize}
        \item \textbf{T1596 (Search Open Technical Databases):} Identificación de subdominios y certificados mediante plataformas públicas como \texttt{crt.sh}.
        \item \textbf{T1589 (Gather Victim Identity Info):} Recolección de correos institucionales de autoridades y docentes desde la web institucional.
    \end{itemize}
    \item \textbf{Acceso Inicial (TA0001):}
    \begin{itemize}
        \item \textbf{T1566.002 (Spearphishing Link):} Envío de enlaces maliciosos simulando comunicaciones del rectorado o la OTI.
        \item \textbf{T1078 (Valid Accounts):} Acceso a plataformas universitarias utilizando credenciales legítimas comprometidas sin factor MFA.
    \end{itemize}
    \item \textbf{Ejecución (TA0002):}
    \begin{itemize}
        \item \textbf{T1059 (Command and Scripting Interpreter):} Ejecución no autorizada de comandos en el intérprete del servidor web a través de vulnerabilidades de aplicación.
    \end{itemize}
    \item \textbf{Persistencia (TA0003):}
    \begin{itemize}
        \item \textbf{T1505.003 (Web Shell):} Carga de archivos maliciosos en directorios de carga (\textit{uploads}) del Campus Virtual o portales web.
    \end{itemize}
    \item \textbf{Escalada de Privilegios (TA0004):}
    \begin{itemize}
        \item \textbf{T1068 (Exploitation for Privilege Escalation):} Aprovechamiento de fallos en versiones desactualizadas de software base (ej. servidores web Nginx o kernels Linux).
    \end{itemize}
    \item \textbf{Descubrimiento (TA0007):}
    \begin{itemize}
        \item \textbf{T1046 (Network Service Discovery):} Escaneo de puertos y servicios internos lanzado desde estaciones de cómputo o laboratorios de alumnos.
    \end{itemize}
    \item \textbf{Movimiento Lateral (TA0008):}
    \begin{itemize}
        \item \textbf{T1021.002 (SMB Shares):} Propagación entre estaciones de trabajo administrativas mediante carpetas compartidas sin segmentar.
    \end{itemize}
    \item \textbf{Exfiltración (TA0010):}
    \begin{itemize}
        \item \textbf{T1048 (Exfiltration Over Alternative Protocol):} Salida de bases de datos o información confidencial mediante canales web cifrados o protocolos no supervisados.
    \end{itemize}
\end{itemize}

\section{Matriz MITRE ATT\&CK Institucional y Brechas de Visibilidad}

\begin{center}
\small
\begin{longtable}{>{\bfseries}p{3.2cm} p{4.0cm} p{4.2cm} p{3.2cm}}
\toprule
\rowcolor{outermarginbgcolor}
Táctica MITRE & Técnica / ID & Escenario en la UNAP & Control / Brecha de Visibilidad \\
\midrule
\endfirsthead
\toprule
\rowcolor{outermarginbgcolor}
Táctica MITRE & Técnica / ID & Escenario en la UNAP & Control / Brecha de Visibilidad \\
\midrule
\endhead
\bottomrule
\endfoot
Reconocimiento & T1596 / T1589 & Rastreo pasivo en crt.sh y directorio docente & Brecha: Información expuesta públicamente \\
\midrule
Acceso Inicial & T1566.002 / T1078 & Phishing con enlaces apócrifos y cuentas sin MFA & Control: Despliegue de MFA y filtrado de correo \\
\midrule
Ejecución & T1059 & Ejecución de scripts en servidores web vulnerables & Control: Restricción de permisos y WAF \\
\midrule
Persistencia & T1505.003 & Carga de Web Shell en formularios sin sanitizar & Brecha: Carencia de monitoreo de integridad FIM \\
\midrule
Escalada Priv. & T1068 & Exploits en Nginx 1.16.1 o Linux legados & Control: Gestión y parcheo de vulnerabilidades \\
\midrule
Descubrimiento & T1046 & Escaneo de red interna desde laboratorios estudiantiles & Brecha: Falta de microsegmentación entre VLANs \\
\midrule
Mov. Lateral & T1021.002 & Movimiento vía SMB en redes de oficinas & Control: Deshabilitar SMBv1 y aislar puestos \\
\midrule
Exfiltración & T1048 & Extracción de tablas de notas vía HTTPS no filtrado & Control: Inspección TLS y DLP perimetral \\
\bottomrule
\end{longtable}
\end{center}

\section{Regla de Detección SIEM (Wazuh / Sysmon)}
Para mitigar la técnica de ejecución y persistencia (T1505.003 - Web Shell), se define la siguiente regla de correlación:

\begin{cmdbox}{Regla Wazuh para Detección de Creación de Web Shell (/etc/rules/local\_rules.xml)}
<group name='web,attack,'>
  <rule id='100250' level='12'>
    <if_group>syscheck</if_group>
    <match>/var/www/html/uploads/</match>
    <description>Alerta Critica: Creacion de archivo ejecutable (.php) en carpeta uploads</description>
    <mitre>
      <id>T1505.003</id>
    </mitre>
  </rule>
</group>
\end{cmdbox}